\section*{Resolución P2}
\textbf{Conjuntos e índices:}
\begin{align*}
J &: \text{Sitios candidatos } (j\in J).\\
K &: \text{Zonas de demanda } (k\in K).\\
T &= \{L,H\}: \text{Tipos de antenas.}
\end{align*}

\textbf{Parámetros:}
\begin{align*}
D_k &: \text{Demanda (número de usuarios) en la zona } k.\\
F_j^t &: \text{Costo fijo de instalación de la antena tipo } t \text{ en } j.\\
C_j^t &: \text{Capacidad nominal de la antena tipo } t \text{ en el sitio } j.\\
B &: \text{Presupuesto total.}\\
u_k &: \text{Utilidad por atender completamente la zona } k.\\
R &: \text{Mínimo número de zonas bien cubiertas.}\\
\alpha &: \text{Fracción mínima de la demanda } D_k \text{ requerida por antena para la condición de redundancia.}\\
a_{kj}^t &\in \{0,1\}: \text{Indicador de cobertura (1 si la antena } t \text{ en } j \text{ cubre la zona } k).\\
M &: \text{Constante Big-M.}
\end{align*}

\textbf{Variables de decisión:}
\begin{align*}
x_{jt} &\in \{0,1\}: \text{ 1 si se instala antena tipo } t \text{ en } j.\\
z_k &\in \{0,1\}: \text{ 1 si la zona } k \text{ es atendida completamente.}\\
w_k &\in \{0,1\}: \text{ 1 si la zona } k \text{ está bien cubierta.}\\
v_{kj} &\in \{0,1\}: \text{ 1 si } j \text{ cubre más de la fracción } \alpha \text{ de la demanda.}\\
d_{kj}^t &\geq 0: \text{ Cantidad de usuarios asignados a la antena tipo } t \text{ en el sitio } j.
\end{align*}

\textbf{Función objetivo:}
\begin{align}
\max \sum_{k\in K} u_k z_k
\end{align}

\textbf{Restricciones:}

\begin{itemize}
    \item Instalación única por sitio
    \begin{equation}
        \sum_{t\in T} x_{jt} \leq 1 \quad  \forall j\in J
    \end{equation}
    \item Cobertura
    \begin{equation}
        d_{kj}^t \leq M a_{kj}^t \quad \forall k\in K,\ \forall j\in J,\ \forall t\in T
    \end{equation}
    \item Capacidad por antena instalada
    \begin{equation}
        \sum_{k\in K} d_{kj}^t \leq C_j^t x_{jt} \quad \forall j\in J,\ \forall t\in T
    \end{equation}
    \item Restricción presupuestaria
    \begin{equation}
        \sum_{j\in J}\sum_{t\in T} F_j^t x_{jt} \leq B
    \end{equation}
    \item Satisfacción de demanda
    \begin{equation}
        \sum_{j\in J}\sum_{t\in T} d_{kj}^t \leq D_k \quad \forall k\in K
    \end{equation}
    \item Atención completa
    \begin{equation}
        \sum_{j\in J}\sum_{t\in T} d_{kj}^t \geq D_k z_k \quad \forall k\in K
    \end{equation}
    \item Condición de bien cubierta
    \begin{itemize}
        \item[a)] Condición de redundancia
        \begin{equation}
            \sum_{j\in J} v_{kj} \geq 2w_k \quad \forall k\in K
        \end{equation}
        \item[b)] Condición de atención completa
        \begin{equation}
            w_k \leq z_k \quad \forall k\in K
        \end{equation}
        \item[c)] Condición de fracción mínima
        \begin{equation}
            \sum_{t\in T} d_{kj}^t \geq \alpha D_k v_{kj} \quad  \forall k\in K,\ \forall j\in J
        \end{equation}
     \end{itemize}
    \item Requisito R
    \begin{equation}
        \sum_{k\in K} w_k \geq R
    \end{equation}
    \item Dominio de las variables
    \begin{equation}
        x_{jt},\,z_k,\,w_k,\,v_{kj} \in \{0,1\}, \quad d_{kj}^t \geq 0.
    \end{equation}
\end{itemize}


Por otro lado, al aumentar el valor de $R$ el problema se hace más restringido, con lo que el espacio factible del problema se reduce. Con esto, el óptimo será peor o igual que el problema original. Como es un problema de maximización, el nuevo valor óptimo será menor o igual al del problema original, pudiendo incluso resultar infactible.